\subsection{Actual transverse line equations and their genuine height bound}
\label{subsec:actual-affine-line-equations}

\begin{definition}[Original parameter polynomial and transverse equations]
\label{algebra:line-parameter-equations}
\lean{EnvelopeProperty.AffineLine.parameterPolynomial,
EnvelopeProperty.AffineLine.lineRelation,
EnvelopeProperty.AffineLine.lineIdeal,
EnvelopeProperty.AffineLine.subst_parameterPolynomial,
EnvelopeProperty.AffineLine.subst_lineRelation,
EnvelopeProperty.AffineLine.lineRelation_self,
EnvelopeProperty.AffineLine.lineRelation_mem,
EnvelopeProperty.AffineLine.lineIdeal_le_ker_subst}
\leanok
\uses{algebra:affine-line-substitution}
Let $F$ be any field, $a,b$ actual coordinate vectors on any index set,
and choose an actual index $j$ with $b_j\ne0$. Define the original
parameter polynomial and transverse equations by
\[
 T_j=b_j^{-1}(X_j-a_j),\qquad
 L_i=X_i-a_i-b_iT_j,\qquad
 J_b=(L_i:i\ne j).
\]
The genuine line substitution sends $T_j$ to the actual parameter $T$
and every original $L_i$ to zero. Its selected equation $L_j$ is
already the zero polynomial, and every equation belongs to $J_b$.
Thus the original transverse ideal lies in the whole substitution kernel.
\end{definition}
\begin{proof}
Apply the actual substitution to the entire parameter polynomial. Its
coefficient factor becomes $b_j^{-1}b_j=1$, yielding the actual polynomial
$T$. Multiplicativity and subtraction then evaluate each entire $L_i$
to $(a_i+b_iT)-a_i-b_iT=0$. In the original polynomial algebra, the
same actual inverse identity makes $L_j=0$. For $i\ne j$, $L_i$ is
an original generator of the ideal; the selected equation belongs by
its proved zero identity. Since every true generator evaluates to zero,
ideal span closure puts the entire original ideal in the whole kernel.
No generator-value assertion is substituted for the resulting kernel containment.
\end{proof}

\begin{theorem}[Genuine whole polynomial quotient identifies the exact line kernel]
\label{algebra:line-quotient-whole-polynomial}
\lean{EnvelopeProperty.AffineLine.lineQuotientHom,
EnvelopeProperty.AffineLine.lineQuotientHom_mk,
EnvelopeProperty.AffineLine.lineQuotientInv,
EnvelopeProperty.AffineLine.lineQuotientHom_comp_inv,
EnvelopeProperty.AffineLine.lineQuotientInv_comp_hom,
EnvelopeProperty.AffineLine.lineQuotientEquiv,
EnvelopeProperty.AffineLine.lineQuotientEquiv_mk,
EnvelopeProperty.AffineLine.ker_subst_eq_lineIdeal}
\leanok
\uses{algebra:line-parameter-equations}
The entire original line substitution induces a genuine $F$-algebra
isomorphism
\[
 F[X_i]/J_b\simeq F[T].
\]
It retains the original substitution on every polynomial. In particular
its whole kernel is exactly the original ideal $J_b$, without radical
replacement. No finite-index, infinite-field or algebraic-closure
hypothesis is needed for this quotient statement.
\end{theorem}
\begin{proof}
The proved actual kernel containment gives a genuine quotient algebra
homomorphism to $F[T]$. Define the opposite homomorphism by sending
its actual parameter $T$ to the original quotient class of $T_j$.
Their composite on $F[T]$ is the identity because the entire parameter
polynomial maps to $T$; the true polynomial algebra universal property
then identifies the whole homomorphism, including all coefficients.

For the opposite composite, use the actual quotient universal property
and the genuine multivariable polynomial algebra universal property.
Every original relation belongs to $J_b$, so its quotient class gives
$[X_i]=[a_i]+b_i[T_j]$ for each actual coordinate, including the selected
one. Original coefficient commutation identifies each scalar class
with its actual scalar image in the quotient. Consequently the
opposite composite is the identity on every original coordinate and
coefficient, hence on every whole polynomial and every quotient class.
These two true inverse identities construct the actual algebra
isomorphism; no quotient-comparison certificate is assumed.

The isomorphism is injective. Thus an original polynomial has zero
substitution image precisely when its original quotient class is zero,
which is precisely membership in the actual original ideal $J_b$.
This proves the entire kernel equality and retains nilpotent ideal
structures in subsequent genuine ideal images.
\end{proof}

\begin{theorem}[Actual finite equation list and real Krull height bound]
\label{algebra:line-equations-height-bound}
\lean{EnvelopeProperty.AffineLine.lineGenerators,
EnvelopeProperty.AffineLine.lineIdeal_eq_span_generators,
EnvelopeProperty.AffineLine.lineIdeal_fg,
EnvelopeProperty.AffineLine.lineGenerators_card_le,
EnvelopeProperty.AffineLine.map_lineIdeal_eq_span_generators,
EnvelopeProperty.AffineLine.height_map_lineIdeal_le}
\leanok
\uses{algebra:line-parameter-equations}
With a finite coordinate index set of cardinality $m$, the original
transverse ideal has a true finite generator list of at most $m-1$
equations. Its entire ideal image under any original polynomial ring
homomorphism $\varphi:F[X_i]\to B$ is generated by the actual images of
that list. If $B$ is noetherian, every actual prime minimal over this
original image ideal has height at most $m-1$.
\end{theorem}
\begin{proof}
Take the actual finite index list, erase the selected parameter
coordinate and map the remaining indices to their original transverse
polynomials. The resulting finite polynomial list has actual underlying
set exactly the original generating set. It therefore spans the whole
original ideal, giving genuine finite generation. Its cardinality is
at most that of the remaining index list, namely $m-1$; possible repeated
or zero equations only reduce the size.

The actual ideal-map/span theorem identifies the entire original ideal
image with the span of the actual polynomial images. Their true finite
list has size at most that of the original list, hence at most $m-1$.
Apply the real noetherian Krull height theorem to a prime genuinely
minimal over this span. Its height is bounded by the number of entries
of the actual list, yielding the stated bound. No component-height or
finite-generation hypothesis encoding this conclusion is assumed.

For a nonzero direction-coordinate coefficient, the preceding actual
kernel theorem identifies this original equation ideal with the line
kernel. The bound is therefore the true $n-1$-equation input for FHX
Proposition~5.22 in $n$ affine coordinates. This precursor does not yet
convert height to quotient dimension, prove every finite-preimage
component dominates the open line, establish geometric integrality
after field extension, or construct its normalized smooth finite-flat
curve. Those are distinct algebraic producers rather than implicit
consequences of this height bound.
\end{proof}
